\documentclass[12pt,a4paper]{article}

\usepackage[utf8]{inputenc}
\usepackage[T1]{fontenc}
\usepackage[brazil]{babel}
\usepackage[top=3cm,left=3cm,bottom=2cm,right=2cm]{geometry}
\usepackage{mathptmx}
\usepackage{inconsolata}
\usepackage{setspace}
\usepackage{indentfirst}
\usepackage{titlesec}
\usepackage{booktabs}
\usepackage{array}
\usepackage{listings}
\usepackage{xcolor}
\usepackage{xurl}
\urlstyle{same}

\onehalfspacing
\setlength{\parindent}{1.25cm}

\titleformat*{\section}{\normalfont\bfseries\normalsize}
\titleformat*{\subsection}{\normalfont\bfseries\normalsize}

\lstset{
    basicstyle=\ttfamily\footnotesize,
    breaklines=true,
    columns=fullflexible,
    keepspaces=true,
    frame=single,
    rulecolor=\color{black!25},
    backgroundcolor=\color{black!4},
    xleftmargin=0.2cm,
    xrightmargin=0.2cm,
    framexleftmargin=4pt,
    framexrightmargin=4pt,
    framextopmargin=4pt,
    framexbottommargin=4pt,
    literate=
        {á}{{\'a}}1 {é}{{\'e}}1 {í}{{\'i}}1 {ó}{{\'o}}1 {ú}{{\'u}}1
        {Á}{{\'A}}1 {É}{{\'E}}1 {Í}{{\'I}}1 {Ó}{{\'O}}1 {Ú}{{\'U}}1
        {ã}{{\~a}}1 {õ}{{\~o}}1 {Ã}{{\~A}}1 {Õ}{{\~O}}1
        {â}{{\^a}}1 {ê}{{\^e}}1 {ô}{{\^o}}1
        {à}{{\`a}}1 {ç}{{\c{c}}}1 {Ç}{{\c{C}}}1
}

\begin{document}

\section{Questão 03: ELIZA, PARRY e ALICE}

\subsection{ELIZA e o roteiro DOCTOR}

Joseph Weizenbaum apresentou o ELIZA em 1966. O programa executava roteiros de conversação. Seu roteiro mais conhecido, DOCTOR, imitava o estilo de um psicoterapeuta rogeriano ao devolver perguntas abertas, reformular partes das falas e convidar o usuário a desenvolver o que havia dito. A separação entre motor e roteiro permitia mudar as regras sem reescrever o programa [1].

\subsection{Regras e algoritmo}

\textbf{1. Seleção de palavras-chave.} O programa examina a entrada e identifica palavras que acionam regras do roteiro. Quando há várias palavras-chave candidatas, sua precedência orienta a ordem de tentativa.

\textbf{2. Decomposição.} Para a palavra-chave selecionada, o motor testa padrões com palavras fixas e segmentos variáveis. Um padrão como (0 EU SINTO 0) encontra a sequência EU SINTO e captura o trecho seguinte. Na notação do artigo, 0 aceita uma quantidade arbitrária de palavras, inclusive nenhuma.

\textbf{3. Remontagem.} Quando um padrão casa com a entrada, a regra de remontagem combina texto fixo com trechos capturados. Substituições de perspectiva ajustam formas como ``meu'' para ``seu''. O roteiro pode alternar entre várias remontagens para um mesmo padrão.

\textbf{4. Memória e resposta de reserva.} O roteiro original podia guardar uma reformulação associada à palavra MY e retomá-la em certas ocasiões. Quando o motor esgota as regras aplicáveis, pode recorrer a uma lembrança disponível nas condições previstas pelo roteiro ou a uma resposta genérica. Esse mecanismo recupera texto preparado anteriormente [1].

\subsection{Diálogo anotado}

O diálogo a seguir foi elaborado em português para demonstrar as regras passo a passo. As falas e os moldes são uma adaptação didática do DOCTOR.

\noindent \textbf{1. Usuário:} Eu sinto medo de errar.\\
\textbf{ELIZA:} Há quanto tempo você sente medo de errar?

A palavra-chave SINTO ativa o padrão (0 EU SINTO 0). O quarto segmento captura ``medo de errar''. A remontagem insere esse trecho em ``Há quanto tempo você sente [4]?''. A pergunta acompanha o tema porque reaproveita as palavras do usuário.

\noindent \textbf{2. Usuário:} Meu pai cobra muito de mim.\\
\textbf{ELIZA:} Conte-me mais sobre seu pai.

Neste roteiro, PAI tem precedência sobre MEU. O padrão (0 MEU PAI 0) identifica a referência familiar. A remontagem usa ``seu pai'' e deixa de fora o trecho sobre a cobrança. A resposta mantém o assunto geral, mas perde a informação que motivou a fala.

\noindent \textbf{3. Usuário:} Porque ele acha que isso ajuda.\\
\textbf{ELIZA:} Essa é a única razão?

PORQUE ativa o padrão (PORQUE 0), que captura ``ele acha que isso ajuda''. A remontagem escolhida devolve uma pergunta fixa e descarta o trecho capturado. A resposta soa plausível no contexto, embora a regra não resolva a quem ``ele'' se refere nem o que ``isso'' significa.

\noindent \textbf{4. Usuário:} Eu sinto que meu computador está lento.\\
\textbf{ELIZA:} Há quanto tempo você sente que seu computador está lento?

O padrão da primeira fala também aceita ``que meu computador está lento''. A remontagem reinsere o trecho com a troca de ``meu'' por ``seu''. O molde trata uma observação sobre um computador como se fosse uma fala sobre sentimentos. Aqui, o roteiro didático reutiliza o mesmo molde para evidenciar a falha.

\noindent \textbf{5. Usuário:} Ontem foi diferente.\\
\textbf{ELIZA:} Pode continuar?

Nenhuma palavra-chave do roteiro didático se aplica e a demonstração usa a resposta genérica de reserva. A pergunta mantém a conversa aberta, mas não investiga o que mudou nem retoma uma entidade das falas anteriores.

\subsection{Evolução dos mecanismos de conversação}

\textbf{ELIZA (1966).} A conversa é conduzida pela forma da entrada. Palavras-chave selecionam regras, padrões extraem trechos e moldes transformam esses trechos em respostas. A memória pode recuperar frases preparadas, com alcance limitado [1].

\textbf{PARRY (1971--1972).} A equipe de Kenneth Colby em Stanford descreveu uma simulação de processos paranoides em 1971 e publicou testes de indistinguibilidade em 1972. PARRY assume o papel de paciente e incorpora um modelo interno que influencia como interpreta perguntas e escolhe respostas. Isso permite sustentar uma postura de desconfiança entre turnos. O RFC 439, publicado em janeiro de 1973, documenta uma conversa entre PARRY e uma implementação do DOCTOR [2--4].

\textbf{ALICE e AIML.} ALICE (Artificial Linguistic Internet Computer Entity), de Richard Wallace, é um chatbot baseado em regras. AIML (Artificial Intelligence Markup Language) é a linguagem usada para escrever categorias de diálogo, normalmente em XML. Cada categoria associa um padrão de entrada a um molde de resposta. Uma categoria pode levar em conta a fala anterior do bot com \texttt{<that>}, o assunto da conversa com \texttt{<topic>} e redirecionar a entrada para outra categoria com \texttt{<srai>}. Esses recursos permitem que ``sim'', por exemplo, acione uma resposta diferente conforme a pergunta anterior [5].

Os três sistemas ilustram escolhas distintas para produzir continuidade. ELIZA explora os próprios termos do usuário. PARRY mantém um modelo interno que orienta o comportamento do personagem. ALICE amplia o repertório de padrões e oferece recursos explícitos de contexto e reutilização de respostas. No diálogo anotado, um mecanismo como \texttt{<that>} poderia ajudar a distinguir respostas curtas a perguntas diferentes. Ainda seria preciso escrever categorias que contemplassem cada situação.

\section*{Referências}

\noindent [1] WEIZENBAUM, J. \textbf{ELIZA: A Computer Program for the Study of Natural Language Communication Between Man and Machine}. Communications of the ACM, 1966. Disponível em: \url{https://web.stanford.edu/class/cs124/p36-weizenabaum.pdf}.

\noindent [2] COLBY, K. M.; WEBER, S.; HILF, F. D. \textbf{Artificial Paranoia}. Artificial Intelligence, 1971. Disponível em: \url{https://doi.org/10.1016/0004-3702(71)90002-6}.

\noindent [3] COLBY, K. M. et al. \textbf{Turing-like indistinguishability tests for the validation of computer simulation of paranoid processes}. Artificial Intelligence, 1972. Disponível em: \url{https://doi.org/10.1016/0004-3702(72)90049-5}.

\noindent [4] CERF, V. \textbf{PARRY Encounters the DOCTOR}. RFC 439, 1973. Disponível em: \url{https://www.rfc-editor.org/rfc/rfc439}.

\noindent [5] WALLACE, R. S. \textbf{AIML 2.0 Working Draft}. ALICE A.I. Foundation, 2014. Disponível em: \url{https://github.com/AIML-Foundation/AIML-2.0-Spec/blob/master/aiml.md}.

\end{document}
